\documentclass[10pt,a4paper]{amsart}
\usepackage[margin=19mm]{geometry}
\usepackage{amsmath,amssymb,mathtools}
\usepackage[scr=rsfs,cal=euler]{mathalfa}
\usepackage{xfrac}
\usepackage[unicode,hidelinks,bookmarks=false]{hyperref}
\newcommand{\ie}{i.e.\ }
\newcommand{\cf}{cf.\ }
\newcommand{\resp}{respectively\ }
\newcommand{\Subsection}{\subsection}
\providecommand{\nobreakdash}{\nobreak\hbox{-}}
\setlength{\emergencystretch}{2em}
\multlinegap0pt
\usepackage{bm}
\usepackage{bbm}
\usepackage{dsfont}
\usepackage{xspace}
\usepackage{cancel}
\usepackage{mathtools}
\usepackage{amsthm}
\makeatletter
\@ifundefined{theorem}{\newtheorem{theorem}{Theorem}}{}
\@ifundefined{lemma}{\newtheorem{lemma}[theorem]{Lemma}}{}
\makeatother
\newcommand{\one}{\mathbf 1}
\title[Settling the Optimal Exponent Relating Sumsets and Differenc]{Settling the Optimal Exponent Relating Sumsets and Difference Sets}
\author{Haowei Lin, Shanda Li}
\date{Source: arXiv:2607.27199v1 (CC BY 4.0)}
\begin{document}
\maketitle

\noindent\emph{Source and licence.} Settling the Optimal Exponent Relating Sumsets and Difference Sets. Version arXiv:2607.27199v1 (CC BY 4.0), distributed under the Creative Commons Attribution 4.0 International licence, as recorded in the arXiv licence metadata for this exact version and shown on the abstract page https://arxiv.org/abs/2607.27199v1. Full rights notice: licenses/arxiv-2607.27199-CC-BY-4.0.txt.\par\smallskip

\noindent\textit{Editorial scope.} Counted unit: the Lemma whose source label is \texttt{lem:difference-recurrence} in the article (source0.tex lines 244--266, source label \texttt{lem:difference-recurrence}), delivered together with the article's own complete original proof of it (source0.tex lines 268--354, one \texttt{proof} environment of 87 lines). No mathematics is written, repaired or re-derived: statement and proof are the article's own text line for line. Required context retained uncounted: none. Every definition, display and label of those lines is retained unchanged. Omitted from this package and not redistributed: 1--243, 267--267, 355--760. Nothing inside a retained excerpt is omitted or altered: every retained body is delivered exactly as the article's lines are written, so no omission receipt of any kind accompanies this package. Citations: the retained excerpts carry no citation command at all, so no bibliography entry of the article is transcribed and no reference is redistributed. Reference adaptation: none was needed and none was made; every reference inside the delivered unit resolves inside it, so no unresolved reference or citation is produced. Printed slips kept literal: no printed word, symbol, hypothesis or number of the counted statement or of its proof was altered. Layout and class: the article's own class is \texttt{amsart} with packages including amsmath, amsthm, amsfonts, array, booktabs, tabularx, mathabx, graphicx, amssymb, babel, hyperref, cleveref, calc, dsfont; the wrapper is the lane's pinned \texttt{amsart} editorial wrapper at 10pt on A4, which restates the theorem-like environments the retained text needs and adds the article's own macro definitions that the retained text needs (\texttt{\textbackslash one}), with extra wrapper packages (bm, bbm, dsfont, xspace, cancel, mathtools). The delivered numbering is the unit's own. Assets: no retained span contains a figure environment, a graphics-inclusion command, a PDF-page inclusion, a table or a TikZ picture; no image file is copied into this package, the package declares no asset dependency and no image file is redistributed with it. Licence: version arXiv:2607.27199v1 of this article is distributed under the Creative Commons Attribution 4.0 International licence, as recorded in the arXiv licence metadata for this exact version and shown on the abstract page https://arxiv.org/abs/2607.27199v1. A full rights notice is delivered at licenses/arxiv-2607.27199-CC-BY-4.0.txt. Characters: every retained excerpt is pure ASCII, so the delivered engine, XeLaTeX with its default Latin Modern text font, reports no missing character.
\begin{lemma}\label{lem:difference-recurrence}
The sequence $t_j$ satisfies
\[
 t_0=1,\qquad t_1=11,
 \qquad
 t_{j+2}=13t_{j+1}-16t_j\quad(j\ge0).
\]
In particular, if
\[
 \lambda=\frac{13+\sqrt{105}}2,
 \qquad
 \mu=\frac{13-\sqrt{105}}2,
\]
then
\begin{equation}\label{eq:t-closed}
 t_j=
 \frac{\sqrt{105}+9}{2\sqrt{105}}\lambda^j
 +
 \frac{\sqrt{105}-9}{2\sqrt{105}}\mu^j
 <\lambda^j
 \qquad(j\ge1).
\end{equation}
\end{lemma}

\begin{proof}
Every residue modulo $12^j$ has a unique balanced expansion
\[
 \sum_{i=0}^{j-1}\varepsilon_i12^i,
 \qquad
 \varepsilon_i\in E:=\{-6,-5,\ldots,5\}.
\]
The residue represented by such a balanced digit string belongs to
\((Y_j-Y_j)\bmod 12^j\) if and only if there are carries
\[
 c_0=0,
 \qquad
 c_i\in\{-1,0,1\},
\]
such that
\begin{equation}\label{eq:carry-equation}
 \delta_i=\varepsilon_i+c_i-12c_{i+1}\in\Delta
 \qquad(0\le i<j).
\end{equation}
Indeed, summing \eqref{eq:carry-equation} after multiplication by $12^i$
shows that the two digit strings differ by $-c_j12^j$.

After any balanced prefix, the only possible nonempty sets of current carries
are
\[
 X_1=\{0\},
 \qquad
 X_2=\{-1,0\},
 \qquad
 X_3=\{0,1\}.
\]
For a current carry-set $X$ and a balanced digit $\varepsilon$, define
\[
 \Phi_\varepsilon(X)=
 \left\{
 c'\in\{-1,0,1\}:
 \varepsilon+c-12c'\in\Delta
 \text{ for some }c\in X
 \right\}.
\]
A check of the twelve possible values of $\varepsilon$ gives the transition
counts
\begin{equation}\label{eq:transition-table}
\begin{array}{c|ccc}
 &X_1&X_2&X_3\\ \hline
 X_1&5&3&3\\
 X_2&4&5&3\\
 X_3&4&4&4
\end{array}
\end{equation}
where rows are current states and columns are next states.  Thus, with
\[
 M=
 \begin{pmatrix}
 5&3&3\\
 4&5&3\\
 4&4&4
 \end{pmatrix},
 \qquad
 e_1=\begin{pmatrix}1\\0\\0\end{pmatrix},
 \qquad
 \one=\begin{pmatrix}1\\1\\1\end{pmatrix},
\]
we have
\[
 t_j=e_1^{\mathsf T}M^j\one.
\]
A direct multiplication gives
\[
 M^2-13M+16I_3
 =
 \begin{pmatrix}
 0&3&-3\\
 0&0&0\\
 0&-4&4
 \end{pmatrix},
 \qquad
 (M^2-13M+16I_3)\one=0.
\]
Therefore
\[
 t_{j+2}-13t_{j+1}+16t_j=0.
\]
The initial values are $t_0=1$ and $t_1=11$.  Solving the recurrence gives
\eqref{eq:t-closed}.  Both coefficients in \eqref{eq:t-closed} are positive
and sum to $1$, while $0<\mu<\lambda$, so $t_j<\lambda^j$ for $j\ge1$.
\end{proof}

\end{document}
